\title{DeepSeekMath: Pushing the Limits of Mathematical Reasoning in Open Language Models}

\begin{abstract}
The intricate and structured demands of mathematical reasoning present substantial difficulties for language models. In this work, we present DeepSeekMath~7B, a model created by extending the pre-training of DeepSeek-Coder-Base-v1.5 7B with an additional 120B tokens predominantly related to mathematics, obtained from Common Crawl, and supplemented with natural language and programming code content. DeepSeekMath~7B attains a notable 51.7\% on the MATH competition benchmark, accomplishing this without dependence on external tools or ensemble methods, and achieving performance comparable to models such as Gemini-Ultra and GPT-4. Moreover, self-consistency across 64 response samples from DeepSeekMath~7B yields a result of 60.9\% on the MATH benchmark. Two principal components contribute to the model’s mathematical reasoning strength: (1) the extraction of valuable data from public web resources via a carefully crafted data selection pipeline, and (2) the implementation of Group Relative Policy Optimization (GRPO), an adaptation of Proximal Policy Optimization (PPO), which not only augments mathematical reasoning proficiency but also streamlines the memory consumption characteristic of PPO.
\end{abstract}

\section{Introduction}

Large language models (LLM) have transformed methodologies for mathematical reasoning within artificial intelligence, driving remarkable progress on quantitative reasoning benchmarks \citep{MATH} as well as geometry reasoning evaluations \citep{trinh2024solving}. In addition, these systems have shown value in supporting human problem-solvers with challenging mathematical tasks \citep{tao}. Despite these successes, state-of-the-art systems like GPT-4 \citep{gpt4} and Gemini-Ultra \citep{gemini} remain unavailable to the public, while the open-source alternatives currently accessible still lag significantly in terms of capability and performance.

This paper presents DeepSeekMath, a domain-specific language model that demonstrably surpasses existing open-source models in mathematical reasoning and approaches the performance exhibited by GPT-4 on established academic tasks. The core advancement is facilitated by constructing the DeepSeekMath Corpus, an extensive, high-quality pre-training resource containing 120B mathematical tokens. The corpus is assembled from the Common Crawl (CC) by employing a fastText-based classification system \citep{joulin2016fasttext}. Initially, this classifier is trained using positive samples drawn from OpenWebMath \citep{openwebmath}, complemented by a broad assortment of unrelated web pages as negative samples. Following this, the classifier is utilized to extract further positive content from the CC, which is then meticulously curated by human annotators. These enhanced data instances serve to retrain the classifier, thereby improving its accuracy. Assessment of the resulting dataset reveals its exceptional quality: our foundational model, DeepSeekMath-Base 7B, achieves 64.2\% on GSM8K \citep{gsm8k} and 36.2\% on the MATH dataset \citep{MATH}, outperforming Minerva 540B \citep{minerva}. Moreover, since the DeepSeekMath Corpus includes multilingual content, the model demonstrates measurable gains on Chinese math benchmarks \citep{wei2023cmath,agieval}. We contend that our approach to curating mathematical data lays a useful foundation for further progress in this area and anticipate significant advancements ahead.

DeepSeekMath-Base is built upon DeepSeek-Coder-Base-v1.5 7B \citep{deepseek-coder}, as empirical results suggest that initializing with a code-focused model is superior to beginning with a generic language model. In addition, our findings reveal that mathematical pre-training not only bolsters performance on mathematics-related tasks but also leads to gains on benchmarks like MMLU \citep{mmlu} and BBH \citep{bbh}, reflecting improvements in the model's broader reasoning skills.

Following the pre-training phase, we refine DeepSeekMath-Base with mathematical instruction tuning leveraging chain-of-thought \citep{cot}, program-of-thought \citep{pot,pal}, and tool-integrated reasoning \citep{tora} datasets. The resultant model, DeepSeekMath-Instruct 7B, outperforms all other 7B-scale models in its class and achieves performance on par with the top 70B open-source instruction-tuned models.

Additionally, we present Group Relative Policy Optimization (GRPO), a modified reinforcement learning (RL) algorithm derived from Proximal Policy Optimization (PPO) \citep{schulman2017proximal}. Distinct from PPO, GRPO eliminates the need for a critic model and, instead, utilizes group-based scores to estimate the baseline. This approach considerably reduces the computational and resource requirements during training. By relying exclusively on a subset of English-language instruction tuning data, GRPO delivers marked improvements over the robust DeepSeekMath-Instruct model, evidenced by significant gains in both in-domain benchmarks (GSM8K: 82.9\% $\rightarrow$ 88.2\%, MATH: 46.8\% $\rightarrow$ 51.7\%) and out-of-domain mathematics tasks (such as CMATH: 84.6\% $\rightarrow$ 88.8\%) during the RL stage. We further articulate a unified framework for interpreting various methodologies, including Rejection Sampling Fine-Tuning (RFT) \citep{yuan2023scaling}, Direct Preference Optimization (DPO) \citep{dpo}, PPO, and GRPO. Through this perspective, we observe that all these approaches can be viewed as forms of direct or streamlined RL methods. Our comprehensive experiments investigate key aspects of this paradigm, considering scenarios such as online versus offline training, supervision by outcome versus by process, and the distinction between single-turn and iterative RL. Finally, we clarify how our RL strategy enhances the performance of instruction-tuned models and outline promising future research directions for further optimizing RL within this unified conceptual framework.

\subsection{Contributions}
Our primary contributions are the development of scalable mathematical pre-training techniques and an in-depth study and analysis of reinforcement learning methodologies.

\textbf{Math Pre-Training at Scale}

\begin{itemize}[topsep=0pt]
    \item Our study demonstrates that the openly available Common Crawl dataset holds substantial potential for mathematical research. Through the development of a carefully curated data filtering pipeline, we assemble the DeepSeekMath Corpus, which comprises 120B tokens sourced from web pages selected specifically for mathematical relevance. This dataset surpasses the scale of math-specific web page corpora used in Minerva \citep{minerva} by nearly sevenfold and exceeds that of the recently introduced OpenWebMath \citep{openwebmath} by a factor of nine.

    \item Our foundational model, DeepSeekMath-Base 7B, attains a level of mathematical performance comparable to Minerva 540B \citep{minerva}, suggesting that model parameter count alone does not solely determine success in mathematical reasoning. High performance is achievable by a comparatively smaller model when pre-trained on superior quality data.
\end{itemize}

\item We present results from our experiments on mathematical training. Pre-training models on code before engaging in math-specific training enhances their proficiency in solving mathematical tasks, regardless of whether auxiliary tools are involved. This provides some evidence in response to the enduring question: \textit{does code training enhance reasoning abilities?} Our findings indicate that, specifically in the context of mathematical reasoning, it does.

\item Despite the widespread practice of training models on arXiv papers—particularly within mathematics-focused research—we observe no significant performance gains on any of the mathematical benchmarks evaluated in this study.

\textbf{Exploration and Analysis of Reinforcement Learning}

\begin{itemize}[topsep=0pt]
    \item We present Group Relative Policy Optimization (GRPO), a resource-efficient and powerful reinforcement learning approach. GRPO eliminates the need for a critic model by deriving the baseline from aggregated group scores, thereby greatly decreasing training costs relative to Proximal Policy Optimization (PPO).
    \item Our results show that GRPO substantially improves the capabilities of the instruction-tuned model DeepSeekMath-Instruct, relying exclusively on instruction-tuning datasets. Moreover, we find that reinforcement learning with GRPO leads to better generalization on out-of-domain tasks.
    \item We propose a cohesive framework that encompasses various strategies, including RFT, DPO, PPO, and GRPO. Through comprehensive empirical studies—such as online versus offline training, outcome-based versus process-oriented supervision, and single-step versus multi-step reinforcement learning—we systematically analyze the critical factors underlying this unified paradigm.
    \item Utilizing our comprehensive paradigm, we delve into the mechanisms that make reinforcement learning effective and outline several promising avenues for advancing large language models’ reinforcement learning.
\end{itemize}

\subsection{Summary of Evaluations and Metrics}

\begin{itemize}[topsep=0pt]
    \item \textbf{English and Chinese Mathematical Reasoning}: We perform extensive evaluations of our models using both English and Chinese datasets, which encompass mathematical tasks ranging from elementary to university-level difficulty. The English datasets utilized are GSM8K \citep{gsm8k}, MATH \citep{MATH}, SAT \citep{llemma}, OCW Courses \citep{minerva}, and MMLU-STEM \citep{mmlu}. For Chinese, we include MGSM-zh \citep{mgsm}, CMATH \citep{wei2023cmath}, Gaokao-MathCloze \citep{agieval}, and Gaokao-MathQA \citep{agieval}. Our assessment measures both the models' proficiency in crafting complete, standalone textual solutions without external tools and their performance on problems solvable via Python.

When evaluated on English benchmarks, DeepSeekMath-Base demonstrates results on par with the proprietary Minerva 540B \citep{minerva} and outperforms all publicly available base models (including Mistral 7B \citep{mistral} and Llemma-34B \citep{llemma}), irrespective of whether those models have received mathematical pre-training, and frequently achieves this by a notable margin. Significantly, DeepSeekMath-Base achieves even stronger outcomes on Chinese benchmarks, which can be attributed to our decision to diverge from previous works \citep{minerva,llemma} that utilize exclusively English math data for pre-training. Instead, our dataset comprises high-quality mathematical data in multiple languages. Through the integration of mathematical instruction tuning and reinforcement learning, our models DeepSeekMath-Instruct and DeepSeekMath-RL exhibit impressive capabilities, marking the first instance in the open-source sphere of surpassing the 50\% accuracy threshold on the advanced MATH competition dataset.

\item \textbf{Formal Mathematics}: We assess DeepSeekMath-Base on the informal-to-formal theorem proving challenge introduced by \citep{dsp_proof}, utilizing the miniF2F dataset \citep{minif2f} and selecting Isabelle \citep{isabelle} as the proof assistant. The results show that DeepSeekMath-Base achieves notable few-shot autoformalization results.
\item \textbf{Natural Language Understanding, Reasoning, and Code}: In order to thoroughly evaluate the models’ abilities in general comprehension, reasoning, and programming, DeepSeekMath-Base is tested on the Massive Multitask Language Understanding (MMLU) benchmark \citep{mmlu}, which spans 57 multiple-choice tasks across a broad range of fields. Additionally, we use BIG-Bench Hard (BBH) \citep{bbh}, which includes 23 tasks that primarily demand complex, multi-step reasoning, along with HumanEval \citep{codex} and MBPP \citep{mbpp}, both established assessments for coding capabilities. Our findings indicate that mathematical pre-training enhances both language comprehension and reasoning skills.

\section{Math Pre-Training}

\subsection{Data Collection and Decontamination} This section details the methodology employed to build the DeepSeekMath Corpus utilizing data from Common Crawl. As illustrated in Figure \ref{fig:data_collect}, we describe a stepwise pipeline designed to efficiently extract a substantial mathematical dataset from Common Crawl, beginning with an initial seed corpus composed of a limited yet highly curated set of mathematical data. It is important to emphasize that this methodology can be extended to additional domains, such as programming.

Initially, we select OpenWebMath \citep{openwebmath}, which consists of a curated set of high-quality mathematical texts from the web, to serve as our foundational seed dataset. Leveraging this seed data, we build a fastText model \citep{joulin2016fasttext} for the purpose of identifying additional mathematical web documents similar in style and content to OpenWebMath. For training, we randomly sample 500,000 entries from OpenWebMath as positive examples and supplement these with 500,000 negative examples drawn from Common Crawl. Training is performed using an open-source fastText implementation\footnote{\url{https://fasttext.cc}}, specifying a 256-dimensional embedding space, a learning rate of 0.1, word n-grams up to length 3, a minimum frequency threshold of 3 for word occurrences, and 3 epochs. To shrink the original Common Crawl dataset, we utilize techniques based on URL deduplication and near-duplicate content removal, yielding a reduced set of 40B HTML webpages. With the trained fastText model, we search for mathematical web pages within the deduplicated Common Crawl data. Subsequently, to ensure only high-quality content is retained, the retrieved pages are sorted by their fastText-assigned scores, and only those ranking highest are selected. The preserved data volume is determined via pre-training experiments conducted on subsets containing the top 40B, 80B, 120B, and 160B tokens. For this initial cycle, we opt to retain the leading 40B tokens.

Following the initial round of data collection, a substantial number of mathematical web pages remain missing, largely due to insufficient variety in the positive samples used to train the fastText model. To address this limitation, we expand the set of initial mathematical web sources to diversify the seed corpus and enhance the performance of the fastText model. The procedure begins by partitioning the entire Common Crawl dataset into non-overlapping domains, with each domain comprising web pages under the same base URL. For each domain, we compute the proportion of web pages captured during the first collection cycle. Domains in which more than 10\% of web pages are gathered are flagged as mathematics-oriented (for instance, \url{mathoverflow.net}).

We then proceed to manually label URLs within these domains that contain mathematical content (such as \url{mathoverflow.net/questions}). Any web pages associated with these labeled URLs that were not previously collected are incorporated into the seed corpus. This method enables the inclusion of a broader set of positive examples, facilitating the development of a refined fastText model that is better equipped to identify mathematical content in future cycles. After executing four rounds of data collection, a total of 35.5M mathematical web pages are obtained, representing 120B tokens. By the fourth round, approximately 98\% of the mathematical data had already been acquired in the third iteration, prompting us to terminate further data collection.

In order to prevent contamination of benchmarks, we adopt the approach outlined in \cite{deepseek-coder}, excluding web content that features questions or answers from well-known English math benchmarks like GSM8K \citep{gsm8k} and MATH \citep{MATH}, as well as Chinese datasets such as CMATH \citep{wei2023cmath} and AGIEval \citep{agieval}. Specifically, we implement a filtering process that removes any text chunk in our math training dataset containing an exact 10-gram match with any substring from these evaluation benchmarks. For benchmark passages shorter than 10 grams but at least 3 grams in length, exact-match filtering is similarly applied to detect and exclude any compromised web documents.

\subsection{Validating the Quality of the DeepSeekMath~Corpus}

We conduct pre-training experiments to evaluate the DeepSeekMath~Corpus in comparison to other recently introduced corpora designed for mathematical training:

\begin{itemize}[topsep=0pt]
    \item \textbf{MathPile} \citep{mathpile}: This is a multi-source dataset comprising 8.9B tokens, compiled from sources such as textbooks, Wikipedia, ProofWiki, CommonCrawl, StackExchange, and arXiv, with more than 85\% of the data originating from arXiv;
    \item \textbf{OpenWebMath} \citep{openwebmath}: Built by filtering CommonCrawl for mathematical text, this collection consists of 13.6B tokens;
    \item \textbf{Proof-Pile-2} \citep{llemma}: This corpus integrates OpenWebMath, AlgebraicStack (which includes 10.3B tokens of mathematical code), and arXiv articles (28.0B tokens). Our experiments with Proof-Pile-2 replicate the arXiv:Web:Code distribution of 2:4:1 as outlined in \cite{llemma}.
\end{itemize}

\subsubsection{Training Setting}
\label{sec:quality-policy}
We fine-tune a general-purpose language model with 1.3 billion parameters—utilizing the same architectural framework as DeepSeek LLMs \citep{deepseek-llm}—referred to as DeepSeek-LLM 1.3B, on various mathematical datasets. Each mathematical dataset is used to train a separate model instance, with a total training volume of 150 billion tokens per dataset. The experiments make use of the resource-efficient HAI-LLM training suite \citep{haillm}. In accordance with the established methodology from DeepSeek LLMs, training employs the AdamW optimizer \citep{adamW} configured with $\beta_1=0.9$, $\beta_2=0.95$, and $\mathrm{weight\_decay}=0.1$. The learning rate schedule is multi-staged: it linearly increases to its maximum after 2,000 warm-up steps, reduces to 31.6\% of the maximum by 80\% progress into training, and further drops to 10.0\% of the peak value upon reaching 90\% of the training. The highest learning rate used is 5.3e-4. For each training step, we utilize a batch size corresponding to 4 million tokens and set the context window to 4,000 tokens.

\subsubsection{Evaluation Results}

\textbf{The DeepSeekMath~Corpus is of high quality, covers multilingual mathematical content, and is the largest in size.}

\begin{itemize}[topsep=0pt]
    \item \textbf{High-quality}: 
    We assess downstream capabilities across 8 mathematical benchmark tasks utilizing few-shot chain-of-thought prompting \cite{cot}. Table \ref{tab:corpora_comparison} highlights a pronounced advantage for the model trained on the DeepSeekMath~Corpus. Additionally, Figure \ref{fig:corpora_comparisons} illustrates that this model surpasses the Proof-Pile-2-trained model after 50B tokens (corresponding to a full epoch of Proof-Pile-2), suggesting the average quality of data within the DeepSeekMath Corpus is superior.
    \item \textbf{Multilingual}:
    The DeepSeekMath~Corpus incorporates resources in a variety of languages, with English and Chinese constituting the principal language groups. According to Table \ref{tab:corpora_comparison}, leveraging the DeepSeekMath~Corpus for training improves mathematical reasoning abilities in both English and Chinese. By comparison, existing mathematical datasets, which are overwhelmingly dominated by English, yield limited or even negative effects on Chinese mathematical reasoning performance.
    \item \textbf{Large-scale}: 
    The scale of the DeepSeekMath~Corpus exceeds that of currently available mathematical datasets by several multiples. Figure \ref{fig:corpora_comparisons} demonstrates that DeepSeek-LLM 1.3B, when trained with DeepSeekMath~Corpus, exhibits a significantly sharper improvement trajectory, along with more sustained progress. By contrast, baseline datasets are considerably smaller and must be iterated over multiple times during training, which results in model performance quickly stagnating.
\end{itemize}

\subsection{Training and Evaluating DeepSeekMath-Base 7B}

This section presents DeepSeekMath-Base 7B, a foundational model designed with robust reasoning capabilities, particularly tailored for mathematical tasks. The model's initialization is based on DeepSeek-Coder-Base-v1.5 7B \citep{deepseek-coder}, after which it undergoes training on 500B tokens. The training dataset consists of 56\% DeepSeekMath Corpus, 4\% AlgebraicStack, 10\% arXiv, 20\% Github code, with the final 10\% comprising natural language data sourced from Common Crawl in both English and Chinese. Training procedures follow those detailed in Section \ref{sec:quality-policy}, with the exception that the learning rate is capped at 4.2e-4 and a batch size of 10M tokens is employed.

We present an extensive evaluation of DeepSeekMath-Base 7B’s mathematical proficiency, emphasizing its performance in independently generating complete mathematical solutions, employing tools to address mathematical tasks, and engaging in rigorous theorem proving. In addition to its mathematical competencies, we further outline a broader characterization of the foundational model, encompassing its aptitude in natural language comprehension, reasoning, and coding abilities.

\paragraph{Mathematical Problem Solving with Step-by-Step Reasoning}

We assess the capabilities of DeepSeekMath-Base in mathematical problem solving by employing few-shot chain-of-thought prompting \citep{cot} on eight different English and Chinese benchmarks. These benchmarks span both quantitative reasoning tasks—such as GSM8K \citep{gsm8k}, MATH \citep{MATH}, and CMATH \citep{wei2023cmath}—and multiple-choice questions, including MMLU-STEM \citep{mmlu} and Gaokao-MathQA \citep{agieval}. This evaluation covers a broad spectrum of mathematical disciplines, ranging from basic arithmetic to advanced collegiate subjects.

Table \ref{tab:base_math_cot} demonstrates that DeepSeekMath-Base 7B achieves superior results on all eight benchmarks compared to other open-source base models, such as the commonly adopted Mistral 7B \citep{mistral} and the more recent Llemma 34B \citep{llemma}, which is further trained with math data from Proof-Pile-2 \citep{llemma}. Remarkably, on the challenging MATH dataset, DeepSeekMath-Base not only exceeds the best open-source base model performances by more than 10\% in absolute terms, but also outperforms Minerva 540B \citep{minerva}—a much larger (77 times) closed-source base model based on PaLM \citep{palm} and further tuned on mathematical material.

\paragraph{Mathematical Problem Solving with Tool Use} We assess the effectiveness of program-assisted mathematical reasoning on the GSM8K and MATH datasets, employing few-shot program-of-thought prompting as described by \citet{pot,pal}. In this framework, models are tasked with constructing a Python script for each problem, leveraging libraries like \textit{math} and \textit{sympy} to carry out complex calculations. The output generated from running these scripts is taken as the model’s final answer. As detailed in Table \ref{tab:base_math_pot_proof}, DeepSeekMath-Base 7B achieves superior results compared to the previous leading model, Llemma 34B.

\paragraph{Formal Mathematics}

The automation of formal proofs plays a crucial role in improving the precision and dependability of mathematical proofs, as well as promoting greater efficiency—a topic that has gained growing interest in recent years. In this study, we assess DeepSeekMath-Base 7B using the informal-to-formal proof task described by \citep{dsp_proof}. The objective of this task is to produce a formal proof given an informal problem statement, its corresponding formal expression, and an informal proof. Our experiments utilize miniF2F \citep{minif2f}, which serves as a standard benchmark for formalizing mathematics at the Olympiad level. For each problem, we employ few-shot prompting to synthesize formal proofs in Isabelle. Consistent with the methodology in \cite{dsp_proof}, our approach involves having models draft proof outlines, after which the existing automated theorem prover Sledgehammer \citep{sledgehammer} completes the proofs by addressing omitted steps. Results summarized in Table \ref{tab:base_math_pot_proof} indicate that DeepSeekMath-Base 7B achieves notable results in the domain of automated proof formalization.

\paragraph{Natural Language Understanding, Reasoning, and Code}

We assess the natural language understanding abilities of our models using MMLU \citep{mmlu}, their reasoning skills on BBH \citep{bbh}, and coding proficiency on HumanEval \citep{codex} and MBPP \citep{mbpp}. According to Table \ref{tab:understanding_reasoning_code}, DeepSeekMath-Base 7B demonstrates marked improvements over its predecessor, DeepSeek-Coder-Base-v1.5 \citep{deepseek-coder}, on both the MMLU and BBH benchmarks, highlighting the beneficial effects of specialized mathematical training for both language comprehension and reasoning tasks. Furthermore, the incorporation of code tokens during continued pretraining allows DeepSeekMath-Base 7B to preserve the coding benchmark performance of DeepSeek-Coder-Base-v1.5. Collectively, DeepSeekMath-Base 7B achieves substantial gains compared to the general-purpose Mistral 7B model \citep{mistral} across all three reasoning and coding evaluations.

\section{Supervised Fine-Tuning}

\subsection{SFT Data Curation}

We assemble a mathematical instruction-tuning dataset that encompasses both English and Chinese mathematical problems, representing a range of disciplines and difficulty levels. Each problem is accompanied by a solution articulated in chain-of-thought (CoT) \citep{cot}, program-of-thought (PoT) \citep{pot,pal}, or tool-integrated reasoning style \citep{tora}. The dataset comprises 776,000 training instances.

\begin{itemize}[topsep=0pt]
    \item \textbf{English mathematical datasets}: We provide tool-integrated solution annotations for GSM8K and MATH datasets, and incorporate selected problems from MathInstruct \citep{MathInstruct} and the training partition of Lila-OOD \citep{lila}, for which CoT or PoT solution formats are available. The English dataset spans multiple branches of mathematics, including but not limited to algebra, probability, number theory, calculus, and geometry.
    \item \textbf{Chinese mathematical datasets}: We curate a collection of K-12 level math problems in Chinese that cover 76 different sub-topics, such as linear equations. Solutions are annotated using both chain-of-thought and tool-integrated reasoning methods.
\end{itemize}

\subsection{Training and Evaluating DeepSeekMath-Instruct 7B}

In this section, we present DeepSeekMath-Instruct 7B, a model refined via mathematical instruction tuning, leveraging DeepSeekMath-Base as the foundation. To prepare training data, examples are stochastically combined until they fill a context window of up to 4K tokens. The model is then trained for 500 iterations, utilizing a batch size of 256 and maintaining a fixed learning rate of $5\times10^{-5}$.

Model performance on mathematical tasks is assessed under two scenarios: both in the absence and presence of tool support. Four quantitative reasoning benchmarks in English and Chinese are employed for evaluation. Our results are compared against top-performing contemporaneous models:

\begin{itemize}[topsep=0pt]
    \item \textbf{Closed-source models} comprise: (1) the GPT series, notably GPT-4 \citep{gpt4} and GPT-4 Code Interpreter~\footnote{\url{https://openai.com/blog/chatgpt-plugins\#\#code-interpreter}}, which are currently the most advanced, (2) Gemini Ultra and Pro \citep{gemini}, (3) Inflection-2 \citep{inflection-2}, (4) Grok-1~\footnote{\url{https://x.ai/model-card}}, along with recently introduced models from Chinese enterprises such as (5) Baichuan-3~\footnote{\url{https://www.baichuan-ai.com}}, and (6) the most recent GLM-4~\footnote{\url{https://open.bigmodel.cn/dev/api\#glm-4}} developed as part of the GLM series \citep{glm}.
\end{itemize}

These models are designed for broad applicability and typically have been subjected to various alignment methodologies.

\item \textbf{Open-source models} comprise: general-purpose models such as (1) DeepSeek-LLM-Chat 67B \citep{deepseek-llm}, (2) Qwen 72B \citep{qwen}, (3) SeaLLM-v2 7B \citep{seallm}, and (4) ChatGLM3 6B \citep{chatglm3}. Additionally, there exist models specialized or enhanced for mathematical tasks: (5) InternLM2-Math 20B~\footnote{\url{https://github.com/InternLM/InternLM-Math}}, based on InternLM2, trained on mathematics-related content followed by instruction tuning; (6) Math-Shepherd-Mistral 7B, which leverages PPO training \citep{schulman2017proximal} on the Mistral 7B architecture \citep{mistral} with a process-supervised reward signal; (7) the WizardMath suite \citep{wizardmath}, enhancing mathematical reasoning skills in Mistral 7B and Llama-2 70B \citep{llama2} models via evolve-instruct (a variant of instruction tuning employing AI-generated instructions) and PPO training, predominantly utilizing problems from GSM8K and MATH; (8) MetaMath 70B \citep{metamath}, a version of Llama-2 70B further refined on an expanded GSM8K and MATH corpus; (9) ToRA 34B \cite{tora}, which is a CodeLlama 34B model specifically fine-tuned for integrated mathematical problem-solving with tool use; and (10) MAmmoTH 70B \citep{MathInstruct}, which applies instruction tuning with MathInstruct on the Llama-2 70B base model.

According to Table \ref{tab:sft_rl_math}, in scenarios where tool assistance is not permitted, DeepSeekMath-Instruct 7B excels at step-by-step mathematical reasoning. On the challenging MATH benchmark, the model outperforms all available open-source systems as well as most proprietary solutions—including Inflection-2 and Gemini Pro—by a margin of no less than 9\% in absolute terms. This advantage holds even against competitors with significantly larger parameter counts (such as Qwen 72B) or those tailored with math-specific reinforcement learning approaches (such as WizardMath-v1.1 7B). Although DeepSeekMath-Instruct demonstrates competitive results compared to Chinese proprietary models GLM-4 and Baichuan-3 on the MATH dataset, it does not yet reach the levels achieved by GPT-4 or Gemini Ultra.

Within the evaluation framework that permits the combination of natural language reasoning and programmatic tool application for solving tasks, DeepSeekMath-Instruct 7B achieves close to 60\% accuracy on the MATH benchmark, exceeding the performance of all prior open-source models. Across additional benchmarks, our approach demonstrates competitive results when compared to DeepSeek-LLM-Chat 67B—the previous leading model, despite being only a tenth of its size.

\section{Reinforcement Learning}

\subsection{Group Relative Policy Optimization} Reinforcement learning (RL) has demonstrated considerable efficacy in enhancing the mathematical reasoning performance of LLMs following the Supervised Fine-Tuning (SFT) phase \citep{wang2023math,wizardmath}. This section outlines our novel and efficient RL strategy, termed Group Relative Policy Optimization (GRPO).

\subsubsection{From PPO to GRPO} Proximal Policy Optimization (PPO) \citep{schulman2017proximal} stands as a predominant actor-critic algorithm utilized during the reinforcement learning fine-tuning phase for large language models (LLMs) \citep{ouyang2022training}. PPO improves LLMs by seeking to maximize the following surrogate objective function:

\begin{equation}
\footnotesize
    \mathcal{J}_{PPO}(\theta) = \mathbb{E}{[q \sim P(Q), o \sim \pi_{\theta_{old}}(O|q)]} \frac{1}{|o|} \sum_{t=1}^{|o|} \min \left[ \frac{\pi_\theta(o_{t} | q, o_{<t})}{\pi_{\theta_{old}}(o_{t} | q, o_{<t})} A_{t}, \text{clip} \left( \frac{\pi_\theta(o_{t} | q, o_{<t})}{\pi_{\theta_{old}}(o_{t} | q, o_{<t})}, 1 - \varepsilon, 1 + \varepsilon \right)  A_{t} \right] ,
\end{equation}

In this context, $\pi_{\theta}$ and $\pi_{\theta_{old}}$ represent the present and preceding policy networks, respectively, while $q$ and $o$ denote questions and corresponding responses, sampled from both the question set and trajectories generated by the previous policy $\pi_{\theta_{old}}$. The parameter $\varepsilon$ is a PPO-specific hyper-parameter used to facilitate stable learning by controlling the degree of policy update through clipping. The advantage estimate, $A_t$, is calculated using Generalized Advantage Estimation (GAE) \citep{gae}, which leverages the rewards $\{r_{\ge t}\}$ as well as predictions from a trainable value network $V_{\psi}$. Consequently, PPO necessitates simultaneous optimization of both the value function and the policy. To counteract excessive reward model optimization, a common strategy is to integrate a token-level KL divergence penalty relative to a reference policy into the per-token reward \citep{ouyang2022training} as follows:

\begin{equation}
    r_{t} = r_\varphi(q, o_{\le t}) - \beta \log\frac{\pi_{\theta}(o_{t}|q, o_{<t})}{\pi_{ref}(o_{t}|q, o_{<t})},
\label{eq:PPO-reward}
\end{equation}

Here, $r_\varphi$ stands for the reward function, $\pi_{ref}$ refers to the reference policy—commonly chosen as the original SFT (supervised fine-tuned) model—and $\beta$ quantifies the KL penalty strength.

Since the value function used in PPO generally constitutes a separate model roughly equivalent in scale to the policy model, this introduces significant demands on both memory and computation. Furthermore, in reinforcement learning settings, the value function primarily serves as a baseline to reduce variance when computing the advantage. In the LLM setting, however, the reward model typically provides feedback only for the final token, which can make it more challenging to train a value function that remains precise across all tokens. To overcome these challenges, as illustrated in Figure \ref{fig:grpo}, we introduce Group Relative Policy Optimization (GRPO), which eliminates the necessity for an auxiliary value function approximation as required in PPO. Instead, GRPO uses the mean reward obtained from several outputs sampled in response to the same prompt as its baseline. Specifically, for each question $q$, GRPO generates a set of outputs $\{o_1, o_2, \cdots, o_G\}$ from the prior policy $\pi_{\theta_{old}}$ and updates the policy model by maximizing the following objective:

\begin{equation}
\footnotesize
\begin{split}
    \mathcal{J}_{GRPO}(\theta) &= \mathbb{E}{[q \sim P(Q), \{o_i\}_{i=1}^G \sim \pi_{\theta_{old}}(O|q)]}  \\
    & \frac{1}{G}\sum_{i=1}^G\frac{1}{|o_i|} \sum_{t=1}^{|o_i|} \left\{ \min \left[ \frac{\pi_\theta(o_{i,t} | q, o_{i,<t})}{\pi_{\theta_{old}}(o_{i,t} | q, o_{i,<t})} \hat{A}_{i,t}, \text{clip} \left( \frac{\pi_\theta(o_{i,t} | q, o_{i,<t})}{\pi_{\theta_{old}}(o_{i,t} | q, o_{i,<t})}, 1 - \varepsilon, 1 + \varepsilon \right)  \hat{A}_{i,t} \right] - \beta \mathbb{D}_{KL}\left[\pi_{\theta} || \pi_{ref}\right]\right\} ,
\end{split}
\label{eq:GRPO-obj}
\end{equation}

where $\varepsilon$ and $\beta$ denote hyper-parameters, and $\hat{A}_{i,t}$ represents the advantage, which is computed using the relative rewards among outputs within each group exclusively; the precise method will be described in the following subsections. The group-relative approach employed by GRPO for advantage computation is particularly compatible with the inherently comparative objectives of reward models, as such models are generally trained on data comprising paired comparisons of answers to the same prompt. It should be highlighted that, rather than incorporating a KL penalty term directly into the reward, GRPO imposes regularization by explicitly including the KL divergence between the current policy and a reference policy in the loss function, thus simplifying the determination of $\hat{A}_{i,t}$. Moreover, unlike the KL penalty formulation adopted in (\ref{eq:PPO-reward}), our method estimates the KL divergence using the unbiased estimator proposed by \citep{kl_approx}:

\begin{equation}
\small
    \mathbb{D}_{KL}\left[\pi_{\theta} || \pi_{ref}\right] = \frac{\pi_{ref}(o_{i,t}|q,o_{i,<t})}{\pi_{\theta}(o_{i,t}|q,o_{i,<t})}- \log\frac{\pi_{ref}(o_{i,t}|q,o_{i,<t})}{\pi_{\theta}(o_{i,t}|q,o_{i,<t})} - 1,
\end{equation}

which always yields a non-negative value.

\begin{algorithm}[t]
  \small
  \caption{Iterative Group Relative Policy Optimization}
  \textbf{Input} initial policy model $\pi_{\theta_{\text{init}}}$; reward models $r_\varphi$; task prompts $\mathcal{D}$;
  hyperparameters $\varepsilon$, $\beta$, $\mu$
  \begin{algorithmic}[1]
    \State Initialize policy model: $\pi_\theta \leftarrow \pi_{\theta_{\text{init}}}$
    \For{iteration = 1 to I}
      \State Set reference model: $\pi_{ref} \leftarrow \pi_\theta$
      \For{step = 1 to M}
      \State Draw mini-batch $\mathcal{D}_b$ from $\mathcal{D}$
      \State Set previous policy: $\pi_{\theta_{old}} \leftarrow \pi_\theta$
      \State For each question $q$ in $\mathcal{D}_b$, generate $G$ samples $\{o_i\}_{i=1}^G$ using $\pi_{\theta_{old}}(\cdot \mid q)$
      \State For every output $o_i$, evaluate rewards $\{r_i\}_{i=1}^{G}$ by applying $r_{\varphi}$
      \State Perform group relative advantage calculation to obtain $\hat{A}_{i,t}$ at token position $t$ for each $o_i$
      \For{GRPO iteration = 1 to $\mu$}
        \State Optimize the policy $\pi_{\theta}$ with respect to the GRPO objective in Equation~\ref{eq:GC-GRPO}
      \EndFor
    \EndFor
    \State Update reward model $r_\varphi$ via ongoing training utilizing a replay buffer.
    \EndFor
  \end{algorithmic}
  \textbf{Output} $\pi_\theta$
  \label{alg:iter-grpo}
\end{algorithm}

\subsubsection{Outcome Supervision RL with GRPO} More specifically, given a question $q$, a collection of candidate outputs $\{o_1, o_2, \cdots, o_G\}$ is drawn from the previous policy $\pi_{\theta_{old}}$. Each of these outputs is evaluated by a reward model, resulting in a set of reward values $\mathbf{r}=\{r_1, r_2, \cdots, r_G\}$. These reward scores are then standardized by removing the mean of the group and dividing by the standard deviation within the group. Within the framework of outcome supervision, the standardized reward is assigned at the conclusion of every output $o_i$, and this same normalized reward is used as the advantage $\hat{A}_{i, t}$ for each token in the output sequence, so that $\hat{A}_{i, t} = \widetilde{r}_i = \frac{r_i- {\rm mean}(\mathbf{r})}{{\rm std}(\mathbf{r})}$. Finally, the policy is updated by optimizing the objective described in equation (\ref{eq:GRPO-obj}).

\subsubsection{Process Supervision RL with GRPO} While outcome supervision furnishes feedback solely at the completion of each output, this approach may prove inadequate and inefficient for guiding the policy, particularly in intricate mathematical tasks. In line with the methodology in \cite{wang2023math}, we adopt process supervision, wherein a reward is assigned at the conclusion of every reasoning step. Specifically, given a question $q$ and $G$ sampled completions $\{o_1, o_2, \cdots, o_G\}$, a process reward model evaluates each step within these outputs, resulting in the following set of rewards: $\mathbf{R} = \{ \{r_1^{index(1)},\cdots,r_1^{index(K_1)}\}, \cdots,  \{r_G^{index(1)},\cdots,r_G^{index(K_G)}\} \}$. Here, $index(j)$ denotes the token position marking the end of the $j$-th reasoning step, and $K_i$ reflects the step count in output $i$. To ensure consistency across rewards, we standardize them by subtracting their mean and dividing by their standard deviation: $\widetilde{r}_i^{index(j)} = \frac{r_i^{index(j)} - {\rm mean(\mathbf{R})}}{{\rm std(\mathbf{R})}}$.

Next, for each token, the process supervision framework computes its advantage as the cumulative sum of all subsequent normalized rewards, as follows: $\hat{A}_{i, t} = \sum_{index(j) \ge t} \widetilde{r}_i^{index(j)}$.
Policy optimization is then performed by maximizing the objective described in equation (\ref{eq:GRPO-obj}).

\subsubsection{Iterative RL with GRPO} During the progression of reinforcement learning, the previously trained reward model may become inadequate for guiding the latest policy model. To address this, we investigate the use of iterative RL with GRPO. As detailed in Algorithm \ref{alg:iter-grpo}, the iterative GRPO procedure involves constructing fresh training datasets for the reward model by leveraging samples collected from the policy model. The reward model is then updated by employing a replay strategy that maintains 10\% of historical samples in the training pool. Subsequently, the policy model is designated as the new reference model, and training continues with this updated policy model using the refreshed reward model.

\subsection{Training and Evaluating DeepSeekMath-RL}

We implement reinforcement learning utilizing DeepSeekMath-Instruct 7B. The RL training corpus is composed of chain-of-thought-style questions drawn from the SFT dataset, specifically related to GSM8K and MATH, amounting to approximately 144K items. To assess the effect of RL on benchmarks without RL training exposure, we omit the remaining SFT questions. The reward model training dataset is constructed according to the methodology in \citep{wang2023math}. For the initial reward model, we use DeepSeekMath-Base 7B and set the learning rate to 2e-5. When applying GRPO, the policy model operates with a learning rate of 1e-6 and a KL coefficient of 0.04. Each question prompts the sampling of $64$ responses. The maximum sequence length is fixed at 1024, with a batch size of 1024 during training. The policy model is updated once per exploration iteration. For evaluation, DeepSeekMath-RL 7B is assessed on the same benchmarks as DeepSeekMath-Instruct 7B. In this setup, GSM8K and MATH tasks involving chain-of-thought reasoning are considered in-domain, whereas all other benchmarks are categorized as out-of-domain tasks.

Table \ref{tab:sft_rl_math} presents comparative results for open- and closed-source models employing both chain-of-thought and tool-based reasoning across English and Chinese evaluation sets. Our observations are as follows: (1) DeepSeekMath-RL 7B achieves accuracy rates of 88.2\% on GSM8K and 51.7\% on MATH when using chain-of-thought reasoning, outperforming all open-source models within the 7B–70B parameter range and most closed-source alternatives. (2) Notably, DeepSeekMath-RL 7B is exclusively fine-tuned with chain-of-thought-style instructional data derived from GSM8K and MATH, initialized from DeepSeekMath-Instruct 7B. Even with this limited training domain, it consistently outperforms DeepSeekMath-Instruct 7B across all evaluation measures, underscoring the advantages conferred by reinforcement learning.

\section{Discussion}
This section details insights gained from our experiments on pre-training and reinforcement learning.

\subsection{Lessons Learnt in Pre-Training}

To begin, we present our observations during the pre-training phase. Except where explicitly noted, our approach aligns with the training configurations described in Section \ref{sec:quality-policy}. Additionally, it is important to mention that references to the DeepSeekMath Corpus throughout this section pertain to an 89B-token dataset obtained from the second round of data collection.

\subsubsection{Code Training Benefits Mathematical Reasoning}

An often-cited but as-yet unproven hypothesis posits that exposure to code enhances reasoning capabilities. In this work, we seek to provide some evidence addressing this claim, focusing specifically on the area of mathematics: code training boosts the mathematical reasoning performance of models regardless of whether external tools are used.

To investigate the influence of code training on mathematical reasoning, we conducted experiments utilizing both two-stage and single-stage training paradigms:

\textbf{Two-Stage Training}

\begin{itemize}[topsep=0pt]
    \item \textbf{Code Training for 400B Tokens $\rightarrow$ Math Training for 150B Tokens}: The DeepSeek-LLM 1.3B model is initially trained on a corpus containing 400B code tokens, after which it undergoes an additional phase of training on 150B tokens from mathematical datasets;
    \item \textbf{General Training for 400B Tokens $\rightarrow$ Math Training for 150B Tokens}: For comparison, we conduct an alternative experiment in which the initial 400B tokens are drawn from a general-purpose corpus (curated by DeepSeek-AI) instead of code. This setup allows us to explore whether exposure to code tokens, relative to general tokens, confers superior abilities in mathematical reasoning.
\end{itemize}

\textbf{One-Stage Training}

\begin{itemize}[topsep=0pt]
    \item \textbf{Math Training with 150B Tokens}: We pretrain DeepSeek-LLM 1.3B solely on 150B math-related tokens;
    \item \textbf{Joint Training with 400B Code Tokens and 150B Math Tokens}: Sequential math training after code pretraining harms coding ability. We therefore explore whether simultaneously mixing code and math tokens during a single training phase can enhance mathematical reasoning capabilities while also mitigating catastrophic forgetting.
\end{itemize}

\paragraph{Results}
Table \ref{tab:code-math} and Table \ref{tab:code-math-general-eval} report the downstream task performance across the various training regimes.

Training with code provides substantial advantages for program-aided mathematical reasoning in both two-stage and one-stage training scenarios. Table \ref{tab:code-math} illustrates that, within the two-stage paradigm, exposure to code alone considerably improves performance on GSM8K and MATH datasets when leveraging Python as the problem-solving language. Further application of mathematical training in the subsequent phase leads to additional gains. Notably, in the one-stage training regime, integrating code tokens with math tokens successfully addresses the problem of catastrophic forgetting characteristic of the two-stage process, while also enhancing both coding capabilities (Table \ref{tab:code-math-general-eval}) and program-based mathematical reasoning (Table \ref{tab:code-math}).

Code training yields gains in mathematical reasoning even when tool use is not involved. In a two-stage training framework, engaging in code training during the first stage produces noticeable, if moderate, improvements. Moreover, this approach enhances the effectiveness of later math-focused training, culminating in superior overall outcomes. In contrast, merging code and math tokens within a single-stage training process appears to undermine mathematical reasoning in the absence of tool use. A possible explanation is that DeepSeek-LLM 1.3B, given its relatively modest size, does not possess sufficient capacity to thoroughly integrate code and mathematical information when exposed to both simultaneously.

\subsubsection{ArXiv Papers Seem Ineffective in Improving Mathematical Reasoning}

ArXiv papers are frequently incorporated into math pre-training datasets \citep{minerva,gpt-f,llemma,mathpile}. Yet, systematic investigations into how these papers influence mathematical reasoning capabilities are sparse. Surprisingly, our experimental findings indicate that arXiv papers provide minimal benefit to mathematical reasoning proficiency. We assessed this outcome across models of multiple scales, such as DeepSeek-LLM 1.3B and DeepSeek-Coder-Base-v1.5 7B \citep{deepseek-coder}, leveraging arXiv datasets processed through different methodologies:

\begin{itemize}[topsep=0pt]
    \item \textbf{MathPile} \citep{mathpile}: a corpus containing 8.9B tokens, curated through heuristic filtering and cleaning procedures, with scientific arXiv papers representing more than 85\% of its content;
    \item \textbf{ArXiv-RedPajama} \citep{redpajama}: comprises all available arXiv LaTeX files after stripping out preambles, comments, macros, and bibliographies, amounting to a total of 28.0B tokens.
\end{itemize}

During our experimental evaluations, we independently pre-train DeepSeek-LLM 1.3B on 150B tokens and DeepSeek-Coder-Base-v1.5 7B on 40B tokens from each arXiv dataset. The results indicate that exposure to arXiv papers does not meaningfully enhance mathematical reasoning capabilities. With training exclusively on arXiv data, neither model demonstrates significant gains—indeed, some results even suggest a decline—across a range of mathematical evaluation benchmarks with varying degrees of complexity. The tested benchmarks encompass quantitative reasoning tasks such as GSM8K and MATH (Table \ref{tab:arxiv-cot}), multiple-choice assessments exemplified by MMLU-STEM (Table \ref{tab:arxiv-cot}), as well as formal mathematical problem sets like miniF2F (Table \ref{tab:arxiv_atp}).

Nevertheless, this finding is subject to certain limitations and should be interpreted cautiously. The present work has not addressed:

\begin{itemize}[topsep=0pt]
    \item How arXiv tokens influence particular mathematics-focused tasks that fall outside the current study, for example, tasks involving the transformation of formal statements or proofs into their informal equivalents (informalization of theorems);
    \item The implications of incorporating arXiv tokens alongside additional data modalities;
    \item If the advantages provided by arXiv papers persist or become more prominent when models of greater scale are considered.
\end{itemize}

Consequently, additional investigation is necessary, and we defer this pursuit to forthcoming research.

\subsection{Insights of Reinforcement Learning}
\subsubsection{Towards to a Unified Paradigm} Here, we introduce a comprehensive paradigm that facilitates the analysis of various training approaches, including SFT, RFT, DPO, PPO, GRPO. We also perform experiments to investigate the elements comprising this unified paradigm. In a general form, the gradient of any training method with respect to the parameter $\theta$ can be represented as:

There are three principal elements in this framework: (1) the \textit{Data Source} $\mathcal{D}$, specifying the set of training samples; (2) the \textit{Reward Function} $\pi_{{rf}}$, which supplies the feedback signal used to guide learning; and (3) the \textit{Algorithm} $\mathcal{A}$, which integrates both the dataset and the reward information to compute the gradient coefficient $GC$, thereby modulating the strength of penalization or encouragement for each training instance. Under this generalized formulation, we examine a selection of prototypical algorithms:

\begin{itemize}
    \item \textbf{Supervised Fine-tuning (SFT)}: This approach adapts a pretrained model to specific tasks by training it on curated SFT datasets selected by humans.
    \item \textbf{Rejection Sampling Fine-tuning (RFT)}: RFT builds on SFT by additionally training the model on a subset of its own generated outputs for SFT prompts, retaining only those responses that are validated as correct.
    \item \textbf{Direct Preference Optimization (DPO)}: DPO extends SFT by further tuning the model on augmented samples it generates, employing a pairwise DPO loss function to encode human preferences.
    \item \textbf{Online Rejection Sampling Fine-tuning (Online RFT)}: Unlike standard RFT, Online RFT starts with the SFT-initialized model but iteratively improves the policy using real-time samples drawn from the current policy itself.
    \item \textbf{PPO/GRPO}: PPO/GRPO initializes the policy with the SFT model and enhances it through reinforcement, leveraging outputs generated from the actively updated policy.
\end{itemize}

The constituent elements of these approaches are outlined in Table \ref{tab:data_policy}. For a comprehensive explanation of the derivation procedures, please see Appendix \ref{app:analysis-rl}.

\paragraph{Commentary on Data Source} We categorize data sources into two primary types: online sampling and offline sampling. In online sampling, training data originates from the real-time interactions generated by the currently trained policy model. In contrast, offline sampling relies on data obtained from the outputs of the initial SFT model. Both RFT and DPO utilize offline data, whereas Online RFT and GRPO employ data produced through online sampling.

As illustrated in Figure \ref{fig:rl-analysis}, our results indicate that Online RFT markedly surpasses RFT on both evaluated benchmarks. At the beginning of training, Online RFT performs similarly to RFT; however, as training progresses, it achieves a clear lead, highlighting the benefits of the online approach. This outcome is expected: during the initial training phases, the actor closely mirrors the SFT model, resulting in sampled data with negligible differences. As training advances, the actor’s outputs diverge more substantially, making real-time data sampling increasingly advantageous.

\paragraph{Observation about Gradient Coefficient} The algorithm utilizes the reward signal to determine the gradient coefficient, which is then used to adjust the model parameters. In our experimental setup, the reward function is separated into `Rule' and `Model' components. The `Rule' approach evaluates responses by directly assessing answer correctness, while the `Model' approach employs a trained reward model to assign scores to responses. This reward model is trained using data labeled by the rule-based evaluation. As illustrated in Equations \ref{eq:GC-RFT} and \ref{eq:GC-GRPO}, a significant distinction emerges between GRPO and Online RFT: GRPO uniquely modulates its gradient coefficient according to the reward value received from the reward model. This capability enables GRPO to apply variable levels of reinforcement or penalty based on the magnitude of each response’s reward. By comparison, Online RFT does not exhibit this flexibility; it applies reinforcement to all correct responses at an identical level and does not impose penalties on incorrect ones.

As illustrated in Figure \ref{fig:rl-analysis}, GRPO outperforms online RFT, underscoring the effectiveness of modifying the positive and negative gradient coefficients. Moreover, GRPO+PS achieves better results than GRPO+OS, which demonstrates the advantages of employing fine-grained, step-dependent gradient coefficients. We further investigate iterative RL by performing two rounds of iteration in our experiments. As depicted in Figure \ref{fig:iter-rl}, we observe that iterative RL yields notable performance gains, with the most pronounced improvement occurring during the initial iteration.

\subsubsection{Why RL Works?} In this study, reinforcement learning is applied to a selected portion of instruction tuning datasets, resulting in notable improvements over the original instruction-tuned model. To better understand the effectiveness of reinforcement learning, we assess both Pass@K and Maj@K accuracies for the Instruct and RL models on two evaluation sets. The results, displayed in Figure \ref{fig:maj-pass}, demonstrate that RL brings about improvements in Maj@K accuracy but does not affect Pass@K. This pattern suggests that the performance gains from RL are primarily due to strengthening the model’s output distribution—specifically, \textbf{the enhancement seems to stem from elevating the correct answer's rank within the TopK outputs, rather than fundamentally increasing core model abilities.} Analogously, \citep{wang2023making} highlight a \textbf{misalignment problem} in reasoning tasks present in SFT models, and their work, along with \citep{yuan2023rrhf,song2023preference,wang2023making}, shows that reasoning performance in SFT models can be increased via various preference alignment techniques.

\subsubsection{How to Achieve More Effective RL?}
\label{sec:effec_rl}
Our experiments indicate that RL can be highly effective for mathematical reasoning problems. Furthermore, we introduce a unified framework to interpret several well-known training strategies, treating each as a manifestation of either direct or approximate RL approaches. According to Equation \ref{eq:objective}, this framework highlights three principal elements: Data Source, Algorithm, and Reward Function. We outline prospective avenues for advancing each of these components.

\paragraph{Data Source} The data source forms the foundational input for all training approaches. Within the RL framework, we define the data source as comprising unlabeled questions paired with responses generated by the policy model. For the experiments in this paper, our data is limited to the set of questions derived from the instruction tuning phase, and output samples are produced using straightforward nucleus sampling. We suspect that this restriction may account for our RL pipeline yielding improvements solely in the Maj@K metric. Future work will extend our investigation to out-of-distribution question prompts, leveraging \textbf{advanced sampling (decoding) strategies} such as tree-search techniques \citep{yao2023tree}. Furthermore, recent developments in \textbf{efficient inference techniques} \citep{xia-etal-2023-speculative,leviathan2023fast,kwon2023efficient,xia2024unlocking}, which critically influence the exploration capabilities of policy models, are anticipated to have a substantial impact as well.

\paragraph{Algorithms} Algorithms utilize both the input data and the reward signal to derive gradient coefficients that adjust the model parameters. According to Equation \ref{eq:objective}, current approaches predominantly \textbf{TRUST} the feedback from the reward function to either increase or decrease the likelihood of producing particular tokens. Nevertheless, the reliability of the reward signal cannot always be guaranteed, particularly in highly intricate tasks. For instance, even in datasets such as PRM800K \citep{lightman2023let}, which are curated by skilled annotators, around 20\% of the annotations are still incorrect\footnote{\url{https://github.com/openai/prm800k/issues/12\#issuecomment-1728491852}}. With this in mind, we aim to investigate reinforcement learning algorithms designed to be robust in the face of noisy reward signals. We contend that \textbf{WEAK-TO-STRONG} \citep{burns2023weak} alignment strategies hold the potential to significantly transform existing learning algorithms.

\paragraph{Reward Function} The reward function serves as the primary source of feedback during training. In reinforcement learning, this function is typically represented by a neural reward model. We identify three pivotal research avenues for reward models: 1) \textbf{Improving the generalization capacity of the reward model.} To effectively boost LLMs’ core abilities, it is crucial for the reward model to generalize well to unseen questions and more complex outputs; if not, reinforcement learning may only reinforce existing patterns rather than lead to substantive progress; 2) \textbf{Capturing the reward model's uncertainty.} Accounting for uncertainty may facilitate integration between less reliable reward models and weak-to-strong training strategies; 3) \textbf{Constructing robust, high-resolution process reward models efficiently} to enable more detailed and informative feedback throughout reasoning tasks \citep{lightman2023let,wang2023math}.

\section{Conclusion, Limitation, and Future Work} This work introduces DeepSeekMath, which achieves superior results compared to existing open-source models on the MATH benchmark at the competition level, and its performance is comparable to leading proprietary systems. DeepSeekMath~is built upon DeepSeek-Coder-v1.5 7B and is subjected to continuous training on 500B tokens, including a substantial portion—120B math-related tokens—extracted from Common Crawl. Our comprehensive ablation analysis reveals that web-derived data has considerable value for enriching mathematical datasets, whereas arXiv does not contribute as significantly as anticipated. We propose Group Relative Policy Optimization (GRPO), a modification of Proximal Policy Optimization (PPO), which substantially enhances mathematical reasoning efficiency while reducing memory demands. Empirical findings demonstrate that GRPO remains beneficial even as DeepSeekMath-Instruct 7B attains strong benchmark results. Additionally, we offer a unified perspective on various methodological approaches and outline several promising avenues for advancing reinforcement learning effectiveness.

While DeepSeekMath~demonstrates strong performance on benchmarks focused on quantitative reasoning, its abilities in geometry and theorem-proving lag behind those of closed-source models. For example, during preliminary testing, the model struggled with tasks involving triangles and ellipses, possibly reflecting a selection bias in the data used for both pre-training and fine-tuning phases. Moreover, due to limitations in model size, DeepSeekMath~performs worse than GPT-4 in few-shot settings. Whereas GPT-4 benefits from few-shot examples to boost its outcomes, DeepSeekMath~yields comparable results in both zero-shot and few-shot scenarios. Moving forward, we plan to enhance our data curation pipeline to assemble a higher-quality pre-training dataset. Additionally, we intend to investigate new strategies (see Section \ref{sec:effec_rl}) to achieve more efficient reinforcement learning for large language models.

\end{document}